% tikz-tensors-canvas.tex -- where a block goes in a picture: left to right on
% a row, rows one under another, relation signs between.
%
% A block is anything a layout places as one piece -- a stack, a grid. Blocks
% are placed from a cursor, on a row centered on y = 0. A block that follows
% another with nothing between them stands beside it; \tneq and \tnapprox
% write a sign after everything on the row and move the cursor past it;
% \tnbreak starts a row under everything drawn so far. So a picture of an
% equation has no coordinates in it: its blocks place one another.
%
% A layout begins a block with \tn@block, which sets the cursor and the row,
% and says how far its block reaches with \tn@reach, which is where a sign or
% the next block goes.
\newdimen\tn@cursor \newdimen\tn@right \newdimen\tn@line@y
\newcount\tn@c \newcount\tn@brk \newcount\tn@fresh \newcount\tn@afterrel
\tikzset{every picture/.append style={execute at begin picture={%
  \global\tn@cursor=0pt \global\tn@right=0pt \global\tn@line@y=0pt
  \global\tn@brk=0 \global\tn@fresh=1 \global\tn@afterrel=0 }}}

% The block reaches at least as far right as <#1>.
\def\tn@reach#1{\tn@len\tn@e{#1}\ifdim\tn@e>\tn@right \global\tn@right=\tn@e\fi}

% \tn@block{<space before, beside a block>}{<how far the block hangs below its
% center line, and the space above it>}
% Begin a block: beside the one before it (after <#1> more), or after the sign
% before it, or at the margin of a new row hung <#2> under everything above.
% \tn@cursor is then where the block starts and \tn@line@y its center line.
\def\tn@block#1#2{%
  \ifnum\tn@fresh=0 \ifnum\tn@afterrel=0
    \global\tn@cursor=\tn@right \global\advance\tn@cursor by #1\relax
  \fi\fi
  \global\tn@fresh=0 \global\tn@afterrel=0
  \ifnum\tn@brk=1
    \tn@anc{current bounding box}{south}\tn@d=\tn@ay\relax
    \tn@len\tn@yz{\the\tn@d-(#2)}%
    \global\tn@line@y=\tn@yz \global\tn@brk=0
  \fi}

% \tnbreak
% Start a new row: the next block goes at the left margin, under everything
% drawn so far. Two rows of one picture are two equations, or the same one
% drawn two ways.
\newcommand{\tnbreak}{%
  \global\tn@brk=1 \global\tn@cursor=0pt \global\tn@right=0pt
  \global\tn@fresh=1}

% \tneq[factor=<text>]       an equals sign
% \tnapprox[factor=<text>]   an approximately-equals sign
% After everything drawn so far, centered on the row like the blocks, with
% <factor> after it if there is one (\tneq[factor=$\lambda$]); the next block
% starts after it.
\pgfkeys{/tn/rel/.is family, /tn/rel/factor/.initial=}
\newcommand{\tn@rel}[2]{% sign, keys
  \pgfkeys{/tn/rel/.cd, factor=, #2}\pgfkeysgetvalue{/tn/rel/factor}\tn@factor
  \node[anchor=west, inner sep=0pt] (tn@eq) at
    ($(\the\tn@right,\the\tn@line@y)+(\tn@gap,0)$)
    {$#1$\ifx\tn@factor\@empty\else\hspace{0.3em}\tn@factor\fi};\tn@tracenode{tn@eq}%
  \tn@anc{tn@eq}{east}\tn@d=\tn@ax\relax
  \global\tn@cursor=\tn@d \global\advance\tn@cursor by \tn@gap\relax
  \global\tn@right=\tn@cursor \global\tn@afterrel=1
}
\newcommand{\tneq}[1][]{\tn@rel{=}{#1}}
\newcommand{\tnapprox}[1][]{\tn@rel{\approx}{#1}}
